\MFQTransition{事件时间仿真答案}

\begin{enumerate}[leftmargin=*]
  \item 对同一事件流运行策略 $A,B$，差值方差为
  \[
  \operatorname{Var}(Y_A-Y_B)=\operatorname{Var}(Y_A)+\operatorname{Var}(Y_B)-2\operatorname{Cov}(Y_A,Y_B).
  \]
  共享到达、方向和完成率会让协方差变大，因而直接比较差值比两个独立均值更精确。共同随机数不能消除模型偏差，只减少比较噪声。
  \item 单个事件的现金 PnL 先写成交方向 $q_i$、半点差 $h$ 和中间价变化 $\Delta m_i$，例如买入后持有到下一事件时 $\pi_i=q_i(\Delta m_i-h)$，卖出时 $q_i=-1$ 自动翻转符号。全路径为 $\sum_i\pi_i$；若还持有库存，必须加上期末按明确价格清算的库存价值，证明该计算记录与现金递推一致。
  \item 对每条完整路径记录 PnL、期末库存和最大绝对库存。均值可用路径样本标准误的 t 区间，尾部则给百分位 bootstrap 或直接报告 P5/P95；比较两个控制器时优先用成对差区间。必须写路径数、随机种子、重采样单位和是否存在共同随机数。
  \item 先选分段常数强度 $\lambda(t)=\lambda_k$，用时间变换 $z_i=\int_{t_{i-1}}^{t_i}\lambda(u)du$。若模型正确，$z_i$ 应近似 IID Exp(1)，所以检查均值、方差、QQ 图和残差自相关；若仍有日内模式，说明强度或事件定义遗漏结构。
  \item 解析输入时先检查 trade id 单调、订单号唯一、价格档不交叉、数量非负；任何失败都返回稳定类别并丢弃整条半更新状态。重复事件不能简单去重后继续，因为这会改变事件时间和队列位置；应记录原始行号和拒绝原因。
  \item 单路径 PnL 既没有抽样误差，也不能揭示库存尾部和成交失败。正确报告要固定配置、生成多路径分布、分解交易数/费用/冲击、给出基线与压力场景，并说明参数是否由同一批路径选择。
  \item SEC 聚合 cancel-to-trade 只告诉市场层面的事件数量比，缺少自己的订单到达序号、前方量、隐藏量、撤单时延和撮合规则。不同队列状态可以产生相同聚合比，因此自己的成交概率不可识别；最多把它当外部先验或压力参数。
  \item Capstone 从冻结事件输入开始，依次重放队列、生成成交、更新库存、生成报价，再比较控制器。若静态策略收益更高，也必须同时报告风险和统计不确定性。
\end{enumerate}

\subsection*{基础衔接：独立计算}
只能说明两策略在该仿真机制下没有差别。二者可能都亏损，也可能共享同一个错误成交假设。
